From - Tue Jun  6 18:34:03 2000
Received: from imo-d03.mx.aol.com (imo-d03.mx.aol.com [205.188.157.35])
	by math.usu.edu (8.9.3+Sun/8.9.1) with SMTP id UAA29553
	for <symanzik@sunfs.math.usu.edu>; Mon, 5 Jun 2000 20:40:49 -0600 (MDT)
From: Case570@netscape.net
Received: from Case570@netscape.net
	by imo-d03.mx.aol.com (mail_out_v27.9.) id z.f0.1ffa8 (16241)
	 for <symanzik@sunfs.math.usu.edu>; Mon, 5 Jun 2000 22:42:37 -0400 (EDT)
Message-ID: <f0.1ffa8.266dbf1d@netscape.net>
Date: Mon, 5 Jun 2000 22:42:37 EDT
Subject:  Stats 5810 Homework # 12
To: symanzik@sunfs.math.usu.edu
MIME-Version: 1.0
Content-Type: text/plain; charset="US-ASCII"
Content-Transfer-Encoding: 7bit
X-Mailer: Unknown
Content-Length: 1673
Status: RO
X-Mozilla-Status: 8001
X-Mozilla-Status2: 00000000
X-UIDL: 38baa60a0000038c

Stats 5810
Homework # 12
Glen Erickson

     After reading the abstracts from the several papers handed out in class I decided to review the one titled:  Interactive Tables And Maps---A Glance At EPA`S Cumulative Exposure Project Web Page. I found this article to be most interesting because I can relate to the way a graph can simpifily everything and make it make sense. When a large group of numbers and rows and rows of data are presented it makes little to no impression in my mind as to what all of this means, however if someone can present it in a graph or a table then the once obscure data takes on a whole new picture. Picture is the main work here and the old saying that a picture paints a thousand words (or replaces tons of data) really applies in statistics. The article describes the way maps and tables were used to help all to see at a glance and digest an overwhelming amount of data and make it make sense. The way that graphs and tables help to organize and help us to visualize data that other wise would tak!
e a great deal of time to sift through is amazing to me. I know that for me that my understanding and retention of facts in greatly enhanced when it is displayed in a visual way, either in a graph or some of easily understandable form.
     I would like to see more data displayed like this, I think more people could understand if they can visualize the data. This concept is not new, but with interactive graphics and the amazing things we can do with computers now days there and limitless possibilities in displaying data sets.

----------
Get your own FREE, personal Netscape Webmail account today at http://home.netscape.com/webmail/
